%% SCRATCH -- superseded figure mock-ups. EDITING VIEW ONLY.
%% Inserted by \paperscratch at the very end of content.tex (after the references).
%% submission.tex and twocolumn.tex define \paperscratch as a no-op, so nothing here can
%% reach a submitted or shared PDF. Nothing in the manuscript body may \ref a figure in
%% this section.
%%
%% Why it exists: on 2026-10-24 the draw.io sketches that stood in for Fig 2 and Fig 6 were
%% replaced by the generated sheets carried in the prelim (see the SOURCE comments above
%% each \tcfig in sections/results.tex). The sketches hold panel ideas the sheets do not
%% (the nested cell representation feeding the iBioFoundry Learn--Design--Build--Test loop,
%% the accuracy-versus-|tau| panel, the data-and-model scaling panel), so they are kept
%% here, with their captions as they stood, rather than deleted.
%% Sources (draw.io, embedded-XML SVG; `make figures` re-exports them from these):
%%   figures/scratch-Fig2-gene-interactions-mockup.pdf
%%     <- notes/assets/drawio/scratch-Fig2-gene-interactions-mockup.drawio.svg
%%        (a copy of Fig2-gene-interactions.drawio.svg as it stood on main; the original
%%        stays in place because feat/025-graph-reg-sweep-readout rewrites it as the real
%%        Fig 2 composition)
%%   figures/scratch-Fig6-metabolism-and-metabolic-engineering-mockup.pdf
%%     <- notes/assets/drawio/scratch-Fig6-metabolism-and-metabolic-engineering-mockup.drawio.svg
%%        (renamed from Fig6-metabolism-and-metabolic-engineering.drawio.svg)
%% On 2026-10-26 three supplementary sketches joined them (SCRATCH Figures 3 to 5), replaced in
%% sections/backmatter.tex by figures generated from torchcell sources (the ontology DAG, the query
%% figure and the ontology-growth figure). Sources renamed with a scratch- prefix so the Makefile's
%% DRAWIO_FIGS rule re-exports only these copies:
%%   figures/scratch-ontology_pydantic_hourglass_data_model.pdf
%%     <- notes/assets/drawio/scratch-ontology_pydantic_hourglass_data_model.drawio.png
%%   figures/scratch-neo4j_cell-conversion-deduplication-aggregation.pdf
%%     <- notes/assets/drawio/scratch-neo4j_cell-conversion-deduplication-aggregation.drawio.png
%%   figures/scratch-FigS-dataset-admission-loop.pdf
%%     <- notes/assets/drawio/scratch-FigS-dataset-admission-loop.drawio
\clearpage
\setcounter{figure}{0}
\renewcommand{\thefigure}{\arabic{figure}}
\renewcommand{\figurename}{SCRATCH Figure}% the SI set "Supplementary Figure"; these are neither
\phantomsection
\addcontentsline{toc}{section}{SCRATCH: superseded figure mock-ups (editing view only)}
\begin{center}
{\large\bfseries\color{stred}SCRATCH: superseded figure mock-ups}\\[2pt]
{\small\color{stred}Not part of the manuscript. Editing view only; absent from submission.tex and twocolumn.tex.}
\end{center}
\noindent These draw.io sketches stood in for Fig.~2 and Fig.~6 until 2026-10-24, when the
generated sheets from the prelim replaced them. They are kept so the panel ideas they carry
are not lost: the nested cell representation feeding a Learn--Design--Build--Test loop with
the NSF iBioFoundry, accuracy as a function of $|\tau|$, data and model scaling, and the
graph-regularization ablation (Fig.~2 plan), and a six-panel metabolism and
metabolic-engineering layout (Fig.~6 plan). Captions are reproduced as they stood. The
flux-balance number in the Fig.~2 sketch caption (0.0006) predates the current sheet, which
prints 0.0019. SCRATCH Figures 3 to 5 are former Supplementary Figures (the hourglass data model,
the three-stage knowledge-graph normalization and the dataset-admission loop), replaced on
2026-10-26 by figures generated from the code; they keep the hourglass framing, the normalization
example and the admission gates with the planned agent loop.
\clearpage % keep the [t] floats below this banner, not above it

%% --- Former Fig 2 mock-up, as it stood on main before the 2026-10-24 swap ---
% Panel plan recorded with the sketch:
%   (a) predicted vs actual trigenic interaction; (b) nested cell representation + iBioFoundry
%   DBTL loop; (c) state of the art vs Yeast9 FBA / DCell / DANGO; (d) accuracy vs |tau|;
%   (e) data & model scaling; (f) graph-regularization ablation. Split into two figures if
%   it gets crowded.
\begin{figure*}[t]
\tcfig{figures/scratch-Fig2-gene-interactions-mockup.pdf}{SCRATCH, superseded mock-up of Fig.~2.
TorchCell for gene--gene interaction prediction
and strain design. \textbf{a}, Higher-order genetic interactions are defined by subtracting
lower-order effects: to isolate a $k$-way interaction you must peel off all interactions
from subcollections of the same genes. For a trigenic interaction, the observed
triple-mutant fitness is compared against the fitness expected from the single- and
double-mutant terms,
$\tau_{ijk}=f_{ijk}-f_if_jf_k-\varepsilon_{ij}f_k-\varepsilon_{ik}f_j-\varepsilon_{jk}f_i$,
with digenic interactions $\varepsilon_{ij}=f_{ij}-f_if_j$.
\textbf{b}, The nested cell representation -- genome, gene networks, metabolism, and
environment -- feeds TorchCell and a Learn--Design--Build--Test loop with the NSF
iBioFoundry.
\textbf{c}, Predicted versus measured trigenic gene--gene interaction.
\textbf{d}, State of the art on trigenic interactions: Pearson $r$ for Yeast9 FBA ($0.0006$),
DCell ($0.157$), DANGO ($0.367$), and TorchCell ($0.454$).}{fig:scratch-ggi-mockup}
\end{figure*}
\FloatBarrier

%% --- Former Fig 6 mock-up, as it stood on main before the 2026-10-24 swap ---
\begin{figure*}[t]
\tcfig{figures/scratch-Fig6-metabolism-and-metabolic-engineering-mockup.pdf}{SCRATCH,
superseded mock-up of Fig.~6. Extending TorchCell to metabolism
and metabolic-engineering strain design. \textbf{a}--\textbf{f}.}{fig:scratch-metabolism-mockup}
\end{figure*}
\FloatBarrier

%% --- Former Supplementary Figure (ontology data model, hourglass), as it stood on main before the 2026-10-26 swap ---
\begin{figure*}[t]
\tcfig{figures/scratch-ontology_pydantic_hourglass_data_model.pdf}{SCRATCH, superseded Supplementary Figure, replaced on 2026-10-26 by the layered
directed acyclic graph of the ontology (\suppfig{fig:ontology-dag}). Ontology data model (hourglass). Each experiment
funnels genotype, phenotype, environment, and experiment-specific metadata through a
shared Data Merge Layer that distinguishes records from the same experiment; the layer
can be ignored when merging into larger datasets.}{fig:scratch-data-model}
\end{figure*}
\FloatBarrier

%% --- Former Supplementary Figure (knowledge-graph normalization), as it stood on main before the 2026-10-26 swap ---
\begin{figure*}[t]
\tcfig{figures/scratch-neo4j_cell-conversion-deduplication-aggregation.pdf}{SCRATCH, superseded Supplementary Figure, replaced on 2026-10-26 by the query figure, whose
panel b draws the same three operators (\suppfig{fig:queried-dataset}). Knowledge-graph normalization in three
stages. Conversion maps phenotype experiment types (GeneEssentiality, SyntheticLethality)
to a common Fitness type; deduplication collapses records equivalent other than their
phenotype-label values; aggregation combines equivalent-genotype records with different
phenotype labels (Fitness, GeneInteraction) into a single HeteroData instance.}{fig:scratch-conversion-dedup-agg}
\end{figure*}
\FloatBarrier

%% --- Former Supplementary Figure (dataset-admission loop), as it stood on main before the 2026-10-26 swap ---
%% Hand-composed schematic (notes/assets/drawio/scratch-FigS-dataset-admission-loop.drawio,
%% renamed from FigS-dataset-admission-loop.drawio on 2026-10-26) of the
%% dataset-admission code path: torchcell/provenance/schema_impact.py (schema-impact gate),
%% torchcell/knowledge_graphs/kg_manifest.py (manifest admission) and
%% torchcell/knowledge_graphs/incremental_import.py with
%% database/slurm/scripts/gilahyper_increment_kg-slurm_docker.slurm (the increment runner), as run
%% for the Nadal-Ribelles 2025 admission (GilaHyper job 1722, 2026-09-11). Panel c is a plan, not
%% a result; see notes/torchcell.knowledge_graphs.dataset-admission-loop.md.
\begin{figure*}[t]
\tcfig{figures/scratch-FigS-dataset-admission-loop.pdf}{SCRATCH, superseded Supplementary Figure, replaced on 2026-10-26 by the six-panel ontology-growth
figure (\suppfig{fig:ontology-growth}); its admission path, gates and planned agent loop are kept here.
\textbf{Adding a dataset without rebuilding the
others, and growing the ontology under the same gates.} \textbf{a},~The admission path for one
dataset. A Zotero item names the paper; the literature mirror holds the paper, its supplements
and data with a manifest that pins every file by sha256 and records how it was retrieved; a
typed loader and its pydantic schema turn the released tables into genotype~$\times$
environment~$\rightarrow$ phenotype records with every sourced value quoted; the dev LMDB build
is verified at levels L0 to L4 (structural, completeness, numeric, units, cross-source); the
BioCypher adapter and graph schema map the records to nodes and edges; the schema-impact gate
diffs the pydantic schema surface against \texttt{main}; and the manifest admission gate
decides whether the served graph can take an increment. \textbf{b},~What the gates decide. Each
loader's dependency closure is the forward closure over the schema surface of the classes it
imports, so a changed class reaches every loader whose closure contains it (loaders 1 and 2
rebuild; 3 and 4 are untouched) while a new class or adapter method reaches only the new loader.
The store manifest records, per served dataset, its closure fingerprints, the graph node classes
and the adapter methods it was serialized under; an increment is admissible only when all three
are unchanged for every served dataset, in which case the dev LMDB is staged, that dataset's CSVs
are emitted, a uniqueness constraint on the shared node id is created, edges the store already
holds are dropped, and \texttt{neo4j-admin database import incremental} runs against the stopped
store before verification, manifest recording and archiving. Incremental import matches existing
node ids and can neither update nor delete them, so any change to a served class is blocked and
takes the full path: the impacted dev LMDBs are rebuilt and re-verified, every dataset's CSVs are
emitted, a new store is imported in full and the manifest is bootstrapped from it.
\textbf{c},~The planned agent loop (dashed borders are not built). An \texttt{add\_dataset} job
starts from a citation key; retrieval uses the mirror, else the recorded method or a web search,
and keeps a copy so a rebuild never depends on the source; the closure neighbors of b and their
papers form the context; N proposing agents draft the loader, the schema delta and the adapter;
M critics score every proposal and the accepted union becomes the plan; one implementer applies
it in one worktree and one pull request, and a failed check (unit tests, type and lint checks,
L0 to L4, schema impact) returns to Propose with the failure. Token usage of every agent run is
priced into a cost ledger. Several jobs run in parallel; the served store takes one increment at
a time, and each increment passes the gates of a. Composed from the code, not from a published
figure.}{fig:scratch-dataset-admission-loop}
\end{figure*}
\FloatBarrier
